% MANUSCRIPT.tex, the manuscript of the theory paper of papers/wealth_garch.
% Created 23 October 2026 as pass 1 of the plan in TODO.md (item PLAN),
% following papers/MANUSCRIPT_SKELETON.md, with
% papers/subjective_racism/paper_a/skeleton_ClassMarkers.tex as the
% reference. Filled the same day with the theoretical rows of the parent
% MANUSCRIPT.tex (author, 23 October 2026: the theory paper rests only on
% its theoretical novelty). TODO.md maps each row to its parent row.
%
% STRUCTURE THE CHECKER RELIES ON (checks/verify_manuscript.py)
%   Rows are written only through the six row macros below, and proofs only
%   through the mproof environment. IDs are K1, K2, ... (introduction),
%   P1, ... (players), A1, ... (assumptions), M1, ... (model), L1, ... (literature)
%   and C1, ... (conclusions).
%     \krow{ID}{headline}{model rows}{scope note}
%     \prow{ID}{player}{choice faced}{how the model settles it}{in the literature}{rows}
%     \arow{ID}{assumption}{in plain English}{symbols}{model rows}{reason}
%     \mrow{ID}{claim in words}{statement in symbols}{Lean theorems}
%     \lrow{ID}{bib key}{what we rely on}{verbatim quote}{page}{model rows}
%     \crow{ID}{statement}{model or literature rows}
%     \begin{mproof}{model row}{Lean theorems} ... \end{mproof}
%   Lists of rows or theorems are comma-separated. Lean theorem names are the
%   fully qualified names declared under theory/lean/.
%
% RULES THE CHECKER ENFORCES
%   Introduction rows point to model rows only. Player rows and conclusion rows
%   point to model or literature rows only. Assumption rows point to model rows
%   only, and every model row rests on at least one assumption. Every model row
%   names Lean theorems that exist in an audited Lean file and has exactly one
%   proof in Appendix A naming the same theorems. Every literature row's quote
%   appears verbatim in the paper's lit/ record. In every proof, plain English
%   is at most 30% of the text, measured as prose words over prose words plus
%   mathematical tokens. Every section and table of the skeleton is present,
%   found by its label.
\documentclass[11pt]{article}
\usepackage[T1]{fontenc}
\usepackage[utf8]{inputenc}
\usepackage{mathptmx}
\usepackage{microtype}
\usepackage[a4paper,margin=1.8cm]{geometry}
\usepackage{amsmath,amssymb,amsthm}
\usepackage{booktabs}
\usepackage{array}
\usepackage{longtable}
\usepackage{url}
\usepackage[hidelinks]{hyperref}
\usepackage{natbib}
\usepackage[scaled=0.92]{helvet}

\emergencystretch=1.5em
% Characters that appear verbatim in quotations of the lit/ records.
\DeclareUnicodeCharacter{2217}{\ensuremath{^*}}
\DeclareUnicodeCharacter{03C0}{\ensuremath{\pi}}
\DeclareUnicodeCharacter{2264}{\ensuremath{\le}}
\DeclareUnicodeCharacter{2212}{\ensuremath{-}}
\DeclareUnicodeCharacter{2206}{\ensuremath{\Delta}}

\newcolumntype{P}[1]{>{\raggedright\arraybackslash}p{#1}}

\urlstyle{tt}
\makeatletter
\g@addto@macro\UrlBreaks{\do\.\do\_\do\,}
\makeatother
\newcommand{\lean}[1]{\path{#1}}
\newcommand{\rowref}[1]{#1}
\newcommand{\unv}[1]{\textit{Unverified:} #1}
\newcommand{\pend}[1]{\textit{Pending:} #1}

\newcommand{\krow}[4]{#1 & #2 & \rowref{#3} & #4 \\ \addlinespace}
\newcommand{\prow}[6]{#1 & #2 & #3 & #4 & #5 & \rowref{#6} \\ \addlinespace}
\newcommand{\arow}[6]{#1 & #2 & #3 & #4 & \rowref{#5} & #6 \\ \addlinespace}
\newcommand{\mrow}[4]{#1 & #2 & #3 & \lean{#4} & p.~\pageref{proof:#1} \\ \addlinespace}
\newcommand{\lrow}[6]{#1 & \citet{#2} & #3 & ``#4'' & #5 & \rowref{#6} \\ \addlinespace}
\newcommand{\crow}[3]{#1 & #2 & \rowref{#3} \\ \addlinespace}

\newenvironment{mproof}[2]
  {\subsection*{Proof of #1}\label{proof:#1}%
   {\raggedright\noindent\textit{Lean:} \lean{#2}.\par}\smallskip}
  {\par\hfill$\square$\par\medskip}

% Working title, from papers/wealth_garch/MANUSCRIPT.tex. The author sets
% the title.
\title{Asymmetric Capital Buffer Control}
\author{}
\date{Manuscript skeleton, 23 October 2026}

\begin{document}
\maketitle

\section{Introduction}
\label{sec:intro}

\subsection{Question}

% The question is the author's, from the parent MANUSCRIPT.tex.
How does an asymmetric penalty on a capital shortfall shape a bank's
recapitalisation policy and the variance of its capital buffer? The bank
recapitalises by impulse control, paying a fixed issuance cost, and takes risky
exposure under a regulatory cap.

\subsection{Answer}

The penalty moves the thresholds. The asymmetry parameter $\lambda_S$ and the
fixed issuance cost $K$ both raise the dividend barrier (M1). The
recapitalisation trigger rises in $\lambda_S$ and falls in $K$ (M2). Because the
two parameters move the trigger in opposite directions, the Jacobian of the
barrier and the trigger in $(\lambda_S,K)$ is non-singular, so the pair is
locally identified from the two thresholds (M3), under the hypotheses that M4
and Appendix~\ref{app:evidence} list.

The penalty does not move the limits of buffer variance. Where the optimal
exposure binds the cap near the distress boundary and deep in surplus, buffer
variance tends to the deficit limit and the surplus limit, which involve only
the asset and liability shocks and the cap (M10, M11). For every exposure that
complies with the cap, binding or not, buffer variance is bounded by the same
quantities (M12). The asymmetry parameter enters the thresholds and not the
limits of buffer variance.

\footnotesize
\begin{longtable}{@{}P{0.8cm}P{7.0cm}P{1.9cm}P{6.4cm}@{}}
\caption{Headlines.}\label{tab:headlines}\\
\toprule
ID & Headline & Model rows & Scope note \\
\midrule
\endfirsthead
\toprule
ID & Headline & Model rows & Scope note \\
\midrule
\endhead
\krow{K1}{The asymmetry parameter $\lambda_S$ and the fixed issuance cost $K$ are locally identified from the dividend barrier and the recapitalisation trigger, because both raise the barrier and they move the trigger in opposite directions.}{M1, M2, M3, M4}{Local identification from the thresholds, not an estimator. Rests on the hypotheses of Appendix~\ref{app:evidence}, among them the verification theorem, the envelope identities, the signs of $W'(x_L)$ and $N'$, $\Psi(x_L)>0$ and $\rho_L>\mu'(y^*)$. Fails if the two parameters moved the trigger in the same direction (M3). The fixed cost is charged per unit of liabilities (Appendix~\ref{app:scope}). The responses of the thresholds to the fixed cost alone are not claimed new, since \citet{peura2006} and \citet{daihuangkeppo2019} tabulate them, nor is pinning the fixed cost and a second parameter from moments of the thresholds, which \citet{peura2006} do by calibration (L7, L8, L9).}
\krow{K2}{The asymmetry parameter enters the thresholds and not the limits of buffer variance. Where the optimal exposure binds the regulatory cap near both ends, those limits depend only on the asset and liability shocks and the cap, and for every compliant exposure buffer variance is bounded by them.}{M3, M10, M11, M12}{That the optimum binds the cap near both ends is a hypothesis, supported only by the numerical solve of the parent bundle (Appendix~\ref{app:evidence}); the bound M12 needs no binding. The surplus limit needs the liquidity branch to bind deep in surplus, which needs borrowing; \citet{bayraktar2026} take no borrowing as the main banking case, under which it never binds (L2, Appendix~\ref{app:scope}). The constant buffer variance on the solvency branch is not claimed new, since it is the constant log-cushion volatility of capped portfolio insurance \citep{blackperold1992}, L5. The contrast is with that rule, whose slopes are set by preferences (L6), while here they are set by regulation and the liability shock is not hedged.}
\bottomrule
\end{longtable}
\normalsize

\section{Model}
\label{sec:model}

\subsection{Players}

\footnotesize
\begin{longtable}{@{}P{0.7cm}P{2.0cm}P{3.2cm}P{4.4cm}P{3.4cm}P{1.5cm}@{}}
\caption{Players of the model, the choice each faces and how the rows settle it.}\label{tab:players}\\
\toprule
ID & Player & Choice faced & How the model settles it & In the literature & Rows \\
\midrule
\endfirsthead
\toprule
ID & Player & Choice faced & How the model settles it & In the literature & Rows \\
\midrule
\endhead
\prow{P1}{Bank}{Its risky exposure, the dividend barrier, the recapitalisation trigger and the injection target.}{The thresholds are derived, in their responses to $\lambda_S$ and $K$ (M1 to M4) and in the target's independence of the starting state (M7). The exposure is not derived. M11 takes it at the cap as a hypothesis, and M12 bounds every compliant exposure.}{The bank of \citet{peura2006} and of \citet{daihuangkeppo2019}, with a fixed issuance cost, and of \citet{bayraktar2026}, with the cap.}{M1, M2, M3, M4, M7, M11, M12, L1, L7, L9}
\prow{P2}{Shareholders}{None in the model. They receive dividends and supply each issue at the fixed and the proportional issuance cost.}{A rule taken as given. The objective is their value net of the shortfall penalty, and the bank chooses for them (M8). Finite issuance follows from the fixed cost (M6).}{In \citet{peura2006} the owners choose the issue.}{M6, M8, L7}
\prow{P3}{Creditors}{None. The liabilities carry their own shock.}{State. The liability volatility and the return-liability correlation enter the limits and bounds of buffer variance (M10 to M12).}{No paper read has a liability shock unhedged by the safe asset; in \citet{blackperold1992} the reserve asset tracks the liability.}{M10, M11, M12, L5}
\prow{P4}{Supervisor}{None in the model. It sets required capital and the regulatory cap.}{A rule taken as given (A2). The shortfall penalty below the reference stands for its restrictions, a reading no row derives (Appendix~\ref{app:scope}).}{The cap of \citet{bayraktar2026}.}{M8, M10, L1}
\bottomrule
\end{longtable}
\normalsize

\subsection{Assumptions}

\footnotesize
\begin{longtable}{@{}P{0.7cm}P{3.5cm}P{4.2cm}P{2.0cm}P{1.5cm}P{3.3cm}@{}}
\caption{Assumptions of the model.}\label{tab:assumptions}\\
\toprule
ID & Assumption & In plain English & Symbols & Model rows & Reason \\
\midrule
\endfirsthead
\toprule
ID & Assumption & In plain English & Symbols & Model rows & Reason \\
\midrule
\endhead
\arow{A1}{The state is $x\in(1,\infty)$, with diffusion $\sigma^2(x,\pi)=\pi^2\sigma^2x^2+2\pi c\sigma\sigma_Lx(1-x)+\sigma_L^2(1-x)^2$ under risky exposure $\pi$, and the capital buffer is $z=\log((x-1)/q)$ with $q>0$.}{The bank holds a risky and a riskless asset and owes liabilities that carry their own shock. The state is the assets-to-liabilities ratio, and the capital buffer measures equity against required capital on a log scale.}{$x$, $\pi$, $\sigma$, $\sigma_L$, $c$, $q$, $z$}{M9, M10, M11, M12}{The balance sheet of \citet{bayraktar2026} in ratio form. Appendix~\ref{app:scope}, the balance sheet.}
\arow{A2}{The regulatory cap is $u(x)=\min\{(x-1)/a_1,(x-a_2)/a_3\}$ with $0<a_1<a_3$ and $a_2<1$, and the exposure complies, $0\le\pi(x)\,x\le u(x)$.}{Regulation caps the risky exposure by the smaller of a solvency branch, binding near distress, and a liquidity branch, binding deep in surplus.}{$u$, $a_1$, $a_2$, $a_3$, $\pi$, $x$}{M10, M11, M12}{The cap of \citet{bayraktar2026}. Appendix~\ref{app:scope}, the regulatory cap.}
\arow{A3}{Buffer variance is $\zeta(x)=z'(x)^2\,\sigma^2(x,\pi(x))$.}{Buffer variance is the variance rate of the capital buffer, by the change of variable from the state.}{$\zeta$, $z$, $x$, $\pi$}{M10, M11, M12}{Itô's formula, not proved in Lean. Appendix~\ref{app:evidence}.}
\arow{A4}{The objective is dividends net of issuance costs minus the running shortfall penalty $\Lambda(x)=(\lambda_S-1)(R-x)^+$ with $\lambda_S\ge1$, discounted at $\rho_L$.}{Shareholders value dividends net of what they pay for issues, less a penalty for each moment the state sits below the reference, weighted by the asymmetry parameter.}{$\Lambda$, $\lambda_S$, $R$, $\rho_L$, $x$}{M1, M2, M3, M4, M8}{The asymmetric penalty is the paper's question. Appendix~\ref{app:scope}, the shortfall penalty.}
\arow{A5}{Each issue costs $K$ per unit of liabilities plus $\kappa$ per unit raised.}{Every recapitalisation pays a fixed issuance cost that does not scale with its size and a proportional issuance cost.}{$K$, $\kappa$}{M2, M5, M6, M7}{The cost structure of \citet{peura2006} and \citet{bolton2011}. Appendix~\ref{app:scope}, the fixed issuance cost.}
\arow{A6}{$V$ is the value function, solves the quasi-variational inequality with impulse operator $\mathcal MV$, is $C^2$ on the inaction region, and the optimal policy is a dividend barrier $y^*$, a trigger $x_L$ and a target $y_{\text{post}}$, with $\partial_{\lambda_S}V=-W$ and $\partial_KV=-N$.}{The value of the bank is that of its best policy, which pays dividends at a barrier and recapitalises from a trigger to a target, and the value responds to the two parameters through the expected discounted shortfall and the discounted injection count.}{$V$, $\mathcal MV$, $y^*$, $x_L$, $y_{\text{post}}$, $W$, $N$, $\lambda_S$, $K$}{M1, M2, M3, M4, M5}{The verification theorem and the envelope identities, not proved in Lean. Appendix~\ref{app:evidence}.}
\bottomrule
\end{longtable}
\normalsize

\subsection{Results}

\footnotesize
\begin{longtable}{@{}P{0.8cm}P{3.8cm}P{6.3cm}P{3.7cm}P{1.1cm}@{}}
\caption{Results of the model, each with its statement and the Lean theorems
that prove it.}\label{tab:results}\\
\toprule
ID & Claim & Statement & Lean & Proof \\
\midrule
\endfirsthead
\toprule
ID & Claim & Statement & Lean & Proof \\
\midrule
\endhead
\mrow{M1}{The dividend barrier rises in the asymmetry parameter and in the fixed issuance cost.}{$\theta\in\{\lambda_S,K\}$, $F=W$ or $N$. If $V''(y^*)=0$, $V'(y^*)=1$, $\Lambda'(y^*)=0$, $F'(y^*)=0$, $\sigma^2(y^*)>0$, $\rho_L\ne\mu'(y^*)$, $\partial_\theta V=-F$ and the chain rule holds along $y^*(\theta)$, then $dy^*/d\theta=\rho_LF(y^*)/(\rho_L-\mu'(y^*))$, positive when $0<\rho_L$, $\mu'(y^*)<\rho_L$ and $F(y^*)>0$.}{Identification.vppp_at_barrier, Identification.fpp_at_barrier, Identification.barrier_slope, Identification.barrier_slope_pos, Identification.control_barrier_sign_needs_gap, Envelope.wode_of_envelope, Envelope.node_of_envelope, Envelope.neumann_at_ystar}
\mrow{M2}{The recapitalisation trigger rises in the asymmetry parameter and falls in the fixed issuance cost, the injection target rises in the fixed issuance cost, and the trigger-to-target gap widens in it.}{If $V''(x)\,dx/d\theta+\partial_\theta V'(x)=0$ along a boundary $x(\theta)$, $V''(x)\ne0$ and $\partial_\theta V'=-F'$, then $dx/d\theta=F'(x)/V''(x)$. With $W'(x_L)>0$, $N'(x_L),N'(y_{\text{post}})<0$ and $V''(x_L)>0>V''(y_{\text{post}})$, $dx_L/d\lambda_S>0$, $dx_L/dK<0$, $dy_{\text{post}}/dK>0$ and $d(y_{\text{post}}-x_L)/dK>0$.}{Identification.boundary_slope, Identification.trigger_rises_in_lambda, Identification.trigger_falls_in_K, Identification.target_rises_in_K, Identification.gap_increasing_in_K, Identification.control_trigger_slope_needs_curvature, Envelope.trigger_jump}
\mrow{M3}{The Jacobian of the barrier and the trigger in the two parameters is non-singular, because the parameters move the trigger in opposite directions.}{$\det\partial(y^*,x_L)/\partial(\lambda_S,K)=\frac{\rho_L}{(\rho_L-\mu'(y^*))V''(x_L)}\big(W(y^*)N'(x_L)-N(y^*)W'(x_L)\big)$, negative under the signs of M1 and M2. For $a_{11},a_{12},a_{21}>0>a_{22}$, $a_{11}a_{22}-a_{12}a_{21}<0$. With $a_{22}>0$ it can vanish.}{Identification.det_formula, Identification.det_neg_of_opposite, Identification.det_neg, Identification.control_same_direction_can_fail}
\mrow{M4}{The value function is convex just above the recapitalisation trigger, strictly so exactly when waiting is strictly suboptimal just below it.}{$\Psi=\rho_LV+\Lambda-(1+\kappa)\mu$. If $\sigma^2(x_L)>0$, $\tfrac12\sigma^2V''+\mu V'-\rho_LV-\Lambda=0$ at $x_L^+$, $V'(x_L)=1+\kappa$ and $\mu(1+\kappa)-\rho_LV-\Lambda\le0$ at $x_L$, then $\tfrac12\sigma^2V''(x_L^+)=\Psi(x_L)\ge0$, and $V''(x_L^+)>0\iff\Psi(x_L)>0$. With $\Psi(x_L)=0$, $V''(x_L^+)=0$.}{SecondOrder.psi_identity, SecondOrder.psi_nonneg_of_qvi, SecondOrder.trigger_convex, SecondOrder.trigger_strict_iff, SecondOrder.control_trigger_flat_when_psi_zero}
\mrow{M5}{The value function has no kink at the recapitalisation trigger.}{If $V\ge\mathcal MV$, $V(x_L)=\mathcal MV(x_L)$, $\mathcal MV$ is affine with slope $1+\kappa$, $V$ has right derivative $V'(x_L^+)$, $\sigma^2(x_L)>0$, and $V'(x_L^+)>1+\kappa$ would make $\tfrac12\sigma^2q+(1+\kappa)\mu-\rho_LV-\Lambda\le0$ hold at $x_L$ for every $q$, then $V'(x_L^+)=1+\kappa$. If $\tfrac12\sigma^2q+(1+\kappa)\mu-\rho_LV-\Lambda\le0$ is needed only for $q\le0$, it holds when $(1+\kappa)\mu\le\rho_LV+\Lambda$. Without $V\ge\mathcal MV$ a downward kink is possible.}{SmoothFit.hasDeriv_Mv, SmoothFit.kink_is_convex, SmoothFit.no_convex_kink, SmoothFit.smooth_fit, SmoothFit.control_kink_needs_obstacle, Envelope.smf_permits}
\mrow{M6}{With a positive fixed issuance cost and a finite value, the discounted number of recapitalisations is finite, so the impulse terms of the verification argument converge absolutely.}{$K>0$, $e_i\ge0$, $K\sum_{i<n}e_i\le B$ for all $n$ $\Rightarrow$ $\sum_{i<n}e_i\le B/K$, $N=\sum_ie_i$ converges and $N\le B/K$. If also $|c_i|\le Ce_i$, $\sum_ic_i$ converges absolutely. With $K=0$, $e_i=1$ the cost bound holds and $N$ diverges.}{ImpulseCount.partial_sums_bounded, ImpulseCount.count_summable, ImpulseCount.count_le, ImpulseCount.impulse_sum_summable, ImpulseCount.control_count_needs_fixed_cost}
\mrow{M7}{With a fixed issuance cost and a linear proportional cost, the injection target does not depend on the state from which the bank injects.}{$\forall z\ \big(V(z)-(1+\kappa)z\le V(y)-(1+\kappa)y\big)\iff\forall z\ \big(V(z)-(1+\kappa)(z-x)-K\le V(y)-(1+\kappa)(y-x)-K\big)$, for every $x$. A target optimal from one $x$ is optimal from every $x'$. With the convex cost $(z-x)^2$ and $V(z)=2z$, the optimal target moves with $x$.}{Targeting.target_free_of_start, Targeting.target_same_for_all_starts, Targeting.control_target_moves_with_convex_cost}
\mrow{M8}{At $\lambda_S=1$ the shortfall penalty vanishes identically, and for every $\lambda_S\ge1$ the running payoff is concave.}{$\Lambda(x)=(\lambda_S-1)(R-x)^+$ $\Rightarrow$ $\Lambda\equiv0$ at $\lambda_S=1$, and $-\Lambda$ is concave on $\mathbb R$ for $\lambda_S\ge1$. At $\lambda_S=0$, $R=0$ it is not concave.}{QVI_Part1.lambda_asymmetry_vanishes_at_one, QVI_Part1.neg_Lambda_concave, QVI_Part1.control_concavity_needs_lambda_ge_one}
\mrow{M9}{The capital buffer maps the state space onto the real line, with required capital at zero and the distress boundary at minus infinity.}{$q>0$ $\Rightarrow$ $z$ is strictly increasing on $(1,\infty)$, $1+qe^{z(x)}=x$ for $x>1$, $z(1+qe^{w})=w$, $z(1+q)=0$, $z(x)\to-\infty$ as $x\to1^+$ and $z(x)\to\infty$ as $x\to\infty$.}{CapGeometry.zmap_inverse, CapGeometry.zmap_of_inverse, Coordinate.zmap_strictMonoOn, Coordinate.zmap_at_requirement, Coordinate.zmap_tendsto_atBot, Coordinate.zmap_tendsto_atTop}
\mrow{M10}{The solvency branch of the cap binds exactly up to the crossing point $\bar x$, and there buffer variance is the constant deficit limit $\nu_1^2$.}{$0<a_1<a_3$, $a_2<1$ $\Rightarrow$ $\bar x-1=a_1(1-a_2)/(a_3-a_1)>0$, $u(x)=(x-1)/a_1$ for $x\le\bar x$, $u(x)=(x-a_2)/a_3$ for $x\ge\bar x$. If also $q>0$, $1<x\le\bar x$ and $\pi(x)\,x=u(x)$, then $\zeta(x)=\nu_1^2$.}{CapGeometry.hasDerivAt_zmap, CapGeometry.zvol2_eq, CapGeometry.xbar_sub_one, CapGeometry.one_lt_xbar, CapGeometry.solvency_binds, CapGeometry.liquidity_binds, CapGeometry.zvol2_solvency}
\mrow{M11}{For any exposure policy that binds the cap near an end of the state space, buffer variance tends to the deficit limit $\nu_1^2$ at the distress boundary and to the surplus limit $\nu_3^2$ deep in surplus.}{$q>0$, $0<a_1<a_3$, $a_2<1$. If $\pi(x)\,x=u(x)$ for all $x$ near $1^+$ then $\zeta(x)\to\nu_1^2$ as $x\to1^+$. If $\pi(x)\,x=u(x)$ for all large $x$ then $\zeta(x)\to\nu_3^2$ as $x\to\infty$. Both fail without $a_1<a_3$ or without binding.}{CapGeometry.deficit_limit, CapGeometry.surplus_limit, CapGeometry.variance_ratio_is_cap_geometry, CapGeometry.control_surplus_limit_needs_a1_lt_a3, CapGeometry.control_constancy_needs_binding}
\mrow{M12}{For every exposure policy that complies with the cap, binding or not, buffer variance is at most the larger of $\sigma_L^2$ and the deficit limit $\nu_1^2$ everywhere, and deep in surplus at most the larger of $\sigma_L^2$ and the surplus limit $\nu_3^2$ up to any margin.}{$q>0$, $a_1>0$, $x>1$, $0\le\pi(x)$, $\pi(x)\,x\le u(x)$ $\Rightarrow$ $\zeta(x)\le\max\{\sigma_L^2,\nu_1^2\}$. If $a_3>0$, $a_2<1$ and compliance holds for all large $x$, then for every $\varepsilon>0$, $\zeta(x)\le\max\{\sigma_L^2,\nu_3^2\}+\varepsilon$ for all large $x$. Both fail with $\pi<0$ or with $\pi x>u(x)$.}{Bounds.qv_le_max, Bounds.zvol2_eq_qv, Bounds.nu2_eq_qv, Bounds.qv_zero, Bounds.deficit_bound, Bounds.surplus_bound, Bounds.control_bound_needs_long_exposure, Bounds.control_bound_needs_compliance}
\bottomrule
\end{longtable}
\normalsize

\section{Literature}
\label{sec:literature}

\footnotesize
\begin{longtable}{@{}P{0.8cm}P{2.2cm}P{3.5cm}P{5.5cm}P{1.3cm}P{1.9cm}@{}}
\caption{Papers read, with the passage each row relies on.}\label{tab:literature}\\
\toprule
ID & Paper & What we rely on & Quote & Page & Rows \\
\midrule
\endfirsthead
\toprule
ID & Paper & What we rely on & Quote & Page & Rows \\
\midrule
\endhead
\lrow{L1}{bayraktar2026}{The cap $u$ and the crossing point $\bar x$ of M10 are this paper's eq.~(4) and its switching point. M10 and M11 are statements about its cap. The parallel ends there, since the paper evaluates no limit of the volatility and has no asymmetry parameter.}{the parameter a3 scales the entire liquidity constraint and determines its asymptotic upper bound}{17}{M10, M11}
\lrow{L2}{bayraktar2026}{The authors' main banking case has no borrowing. With no borrowing the liquidity branch lies above the ceiling once $a_2\le(1-a_3)x$, so it never binds (\lean{BCVW.ceiling_tighter_than_lcr}), and M11's surplus limit belongs to their unrestricted benchmark.}{we treat the π ≤ 1 case as the main banking interpretation and the unrestricted case as a mathematical benchmark}{19}{M11}
\lrow{L3}{bayraktar2026}{The paper has no fixed issuance cost, fixes its recapitalisation threshold exogenously, and names a fixed cost as a change that could move recapitalisation. M2 and M3 take up that change.}{If one introduced fixed issuance costs, endogenous market access, or adverse-selection costs that rise near distress, earlier recapitalization or inaction at the boundary could become optimal.}{13}{M2, M3}
\lrow{L4}{bayraktar2026}{A fixed issuance cost charged in dollars breaks the reduction to one dimension. M1 to M3 charge $K$ per unit of liabilities, which keeps it (Appendix~\ref{app:scope}).}{Fixed issuance costs, nonlinear issuance costs, taxes, deposit insurance premia, endogenous deposit rates, balance-sheet size effects, multiple asset classes, or regulatory constraints depending on dollar levels rather than ratios would generally break the one-dimensional reduction.}{9}{M1, M2, M3}
\lrow{L5}{blackperold1992}{Their rule is the regulatory cap with $a_2=0$, multiple $1/a_1$ and leverage limit $1/a_3$ (\lean{BlackPerold.cap_is_cppi}), and its log-cushion volatility on the multiplier branch is constant (their Prop.~2), the constancy of $\nu_1^2$ in M10. The parallel ends at the liability, which their reserve asset tracks and ours does not.}{maintains the portfolio's risk exposure a constant multiple of the excess of wealth over a floor, up to a borrowing limit}{403}{M10, M11}
\lrow{L6}{blackperold1992}{In their Prop.~11 the multiple is set by the impatience to consume and the leverage limit by risk aversion, so the two volatility levels carry preference information. In M11 the slopes are regulatory and the asymmetry parameter is absent from the limits.}{The multiple depends on the utility function only through}{421}{M11}
\lrow{L7}{peura2006}{The bank form of the class, with a fixed issuance cost as a fixed proportion of size, the device of M1 to M3, and each issue ending at the dividend barrier, a single target as in M7.}{As in the general case, the dividend barrier is also the level up to which the capital ratio is replenished each time a new capital issue is implemented.}{2176}{M1, M2, M7}
\lrow{L8}{peura2006}{They pin the issue delay and the fixed cost by calibration from issue cost and frequency, moments of the thresholds, and their Table 2 has the sign pattern of M3 in the delay and the fixed cost at small costs (\lean{PeuraKeppo.table2_det_neg}). The parallel ends at the parameter, a delay rather than a preference, and at the method, a calibration rather than a local identification theorem.}{For D and K, we identify common estimates across banks such that their implied capital raising cost and frequency are of plausible magnitude.}{2184}{M3}
\lrow{L9}{daihuangkeppo2019}{Their Table 2.1 gives the responses of the dividend and recapitalisation barriers to seven parameters, the fixed cost with the signs of M1 and M2 (\lean{DaiHuangKeppo.table21_K_signs}). They state no identification and have no shareholder asymmetry parameter, and their delay and fixed cost move the trigger the same way.}{K 0.0017 −0.0543 0 −0.0004}{76}{M1, M2, M3}
\lrow{L10}{bolton2011}{Its authors state that the double barrier of payout and external financing is an established structure. Its fixed issuance cost is proportional to capital to keep the model homogeneous, as in M1 to M3, and each issue returns to one target, as in M7.}{While this double-barrier policy is not new (it goes back to the inventory model of Miller and Orr 1966 in corporate finance), our model provides substantial new insight}{3}{M1, M7}
\lrow{L11}{lokka2008}{With proportional issuance costs only, issuance is a minimal reflection at zero, neither a lump nor a target. The fixed cost is what produces the single target of M7.}{In this problem, it is optimal to issue minimal new equity so as to reflect the reserves process at 0 and make minimal dividend payments so as to reflect the reserves process at a certain level x∗ (see Section 4).}{2}{M7}
\lrow{L12}{barberis2001}{With a constant level of loss aversion, the volatility of log returns equals that of dividend growth for every level, the parallel to the absence of $\lambda_S$ from the limits of M11. The parallel ends at the mechanism, a constant price-dividend ratio in an endowment economy rather than a binding regulatory cap.}{the volatility of log returns in this model is equal to the volatility of log dividend growth, namely 12 percent.}{47}{M11}
\lrow{L13}{engle2018}{They split observed equity volatility asymmetry into a mechanical-leverage share and a market-exposure share. The parallel ends at the object, a response of variance to the sign of news rather than limits of variance across regions of a state variable.}{Volatility asymmetry is mostly explained by exposure to the aggregate market, not a mechanical leverage effect.}{449}{M11}
\lrow{L14}{li2023}{In an unconstrained loss-averse control the large-wealth limits of the optimal controls depend on the preference parameters, the contrast to M11, where a binding regulatory cap fixes them.}{our asymptotic limits differ significantly from the ones in the Merton's problem, which now sensitively depends on the reference degree parameter}{13}{M11}
\bottomrule
\end{longtable}
\normalsize

Every paper the tables name has been read in full and has a record with Lean
under \path{lit/}.

\section{Conclusions}
\label{sec:conclusions}

Table~\ref{tab:concl} holds the conclusions that rest on model rows and on
papers read in full. Table~\ref{tab:pending} holds every claim that a result
is new and every other claim still waiting, and each of its rows names what
it waits on.

\footnotesize
\begin{longtable}{@{}P{0.8cm}P{12.3cm}P{3.4cm}@{}}
\caption{Conclusions resting on model rows and papers read.}\label{tab:concl}\\
\toprule
ID & Statement & Rows \\
\midrule
\endfirsthead
\toprule
ID & Statement & Rows \\
\midrule
\endhead
\crow{C1}{Under the hypotheses of Appendix~\ref{app:evidence}, the asymmetry parameter $\lambda_S$ and the fixed issuance cost $K$ are locally identified from the dividend barrier and the recapitalisation trigger. Both parameters raise the barrier, and they move the trigger in opposite directions, so the Jacobian of the two thresholds is non-singular.}{M1, M2, M3, M4}
\crow{C2}{Wherever the optimal exposure binds the cap near both ends, the limits of buffer variance are $\nu_1^2$ and $\nu_3^2$, functions of $(\sigma,\sigma_L,c,a_1,a_3)$ only, and for every compliant exposure buffer variance is bounded by $\max\{\sigma_L^2,\nu_1^2\}$. The asymmetry parameter moves the thresholds (C1) and is absent from these limits and bounds. In capped portfolio insurance, whose geometry the cap shares, the slopes come from preferences, so the same limits would carry them.}{M3, M10, M11, M12, L5, L6}
\crow{C3}{The surplus limit $\nu_3^2$ requires the liquidity branch to bind deep in surplus, which requires borrowing to hold the risky asset. With no borrowing, the main banking case of \citet{bayraktar2026}, the liquidity branch never binds once $a_2\le(1-a_3)x$, so the deficit limit and the bound of M12 are the results that apply there.}{M11, M12, L2}
\bottomrule
\end{longtable}

\begin{longtable}{@{}P{0.8cm}P{12.3cm}P{3.4cm}@{}}
\caption{Claims not yet verified, each naming what it waits on.}\label{tab:pending}\\
\toprule
ID & Statement & Rows \\
\midrule
\endfirsthead
\toprule
ID & Statement & Rows \\
\midrule
\endhead
\crow{C4}{The joint local identification of C1, of a shareholder asymmetry parameter and a fixed issuance cost from the two thresholds by the signs of their responses, is new. Its numerical precedents in other parameters are \citet{peura2006} and \citet{daihuangkeppo2019} (L8, L9). \pend{the identification section of Alvarez, Le Bihan and Lippi (2014), Banque de France WP 492; Glover, Johnson, Evatt and Cheng (2023); Lambertini and Uysal (2013, 2014); each read in full with a record (\path{notes/novelty_check_20261023.md}, section 5).}}{M3, L8, L9}
\crow{C5}{That under a regulatory two-branch cap with an unhedged liability shock the limits of buffer variance carry no information about the asymmetry parameter, while the thresholds do (C2), is new. \pend{Choulli, Taksar and Zhou (2003), SIAM J. Control Optim. 41(6), read in full with a record; and a proof that the optimal exposure binds the cap near both ends, the hypothesis of M11.}}{M11, L5, L6}
\bottomrule
\end{longtable}
\normalsize

\bibliographystyle{plainnat}
\bibliography{refs}

\appendix

\section{Proofs}
\label{app:proofs}

\paragraph{Notation.} The state is the assets-to-liabilities ratio
$x\in(1,\infty)$, $q>0$ required capital as a ratio, and
$z(x)=\log\big((x-1)/q\big)$ the capital buffer. The risky exposure is $\pi(x)$,
and the regulatory cap gives the capped exposure
$u(x)=\min\{(x-1)/a_1,(x-a_2)/a_3\}$ with crossing point
$\bar x=(a_3-a_1a_2)/(a_3-a_1)$. The diffusion of $x$ is
$\sigma^2(x,\pi)=\pi^2\sigma^2x^2+2\pi c\sigma\sigma_Lx(1-x)+\sigma_L^2(1-x)^2$,
with $\sigma$ the asset volatility, $\sigma_L$ the liability volatility and $c$
the return-liability correlation. Buffer variance is
$\zeta(x)=z'(x)^2\,\sigma^2\big(x,\pi(x)\big)$, the exposure ratio is
$\omega(x)=\pi(x)\,x/(x-1)$, and
$\nu_i^2=\sigma^2/a_i^2-2c\sigma\sigma_L/a_i+\sigma_L^2$ for $i\in\{1,3\}$.
$V$ is the value function, $y^*$ the dividend barrier, $x_L$ the
recapitalisation trigger and $y_{\text{post}}$ the injection target. $W$ is the
expected discounted shortfall and $N$ the discounted injection count under the
optimal policy, $\theta$ is $\lambda_S$ or $K$, and $F$ is $W$ or $N$
accordingly. $\rho_L$ is the liability-net discount rate, $\mu$ the drift of
$x$ along the cap and $\mu'$ its derivative in $x$, $\Lambda$ the shortfall
penalty and $R$ its reference, and $\ell=(R-x)^+$ for $F=W$ and $\ell=0$ for
$F=N$. Primes are derivatives in $x$. $\kappa$ is the proportional issuance
cost, $\mathcal MV$ the impulse operator, affine with slope $1+\kappa$ on the
intervention side, and $\mu=\mu(x_L)$ in M4 and M5. In M6, $e_i$ are the
discounted impulse times, $c_i$ the impulse terms and $B$ the cost bound.

\begin{mproof}{M1}{Identification.vppp_at_barrier, Identification.fpp_at_barrier, Identification.barrier_slope, Identification.barrier_slope_pos, Identification.control_barrier_sign_needs_gap, Envelope.wode_of_envelope, Envelope.node_of_envelope, Envelope.neumann_at_ystar}
\[ \tfrac12\sigma^2V'''+\tfrac12(\sigma^2)'V''+\mu V''+\mu'V'-\rho_LV'-\Lambda'=0,\quad V''=0,\ V'=1,\ \Lambda'=0\ \Rightarrow\ \tfrac12\sigma^2V'''=\rho_L-\mu' . \]
\[ \tfrac12\sigma^2F''+\mu F'-\rho_LF+\ell=0,\quad F'(y^*)=0,\ \ell(y^*)=0\ \Rightarrow\ \tfrac12\sigma^2F''=\rho_LF . \]
\[ V'''\frac{dy^*}{d\theta}+\partial_\theta V''=0,\quad \partial_\theta V''=-F''\ \Rightarrow\ \frac{dy^*}{d\theta}=\frac{F''}{V'''}=\frac{\rho_LF(y^*)}{\rho_L-\mu'(y^*)} . \]
Control, $\rho_L=1$, $\mu'=2$, $F=1$ $\Rightarrow$ $dy^*/d\theta=-1$.
\end{mproof}

\begin{mproof}{M2}{Identification.boundary_slope, Identification.trigger_rises_in_lambda, Identification.trigger_falls_in_K, Identification.target_rises_in_K, Identification.gap_increasing_in_K, Identification.control_trigger_slope_needs_curvature, Envelope.trigger_jump}
\[ V''(x)\frac{dx}{d\theta}+\partial_\theta V'(x)=0,\quad \partial_\theta V'=-F'\ \Rightarrow\ \frac{dx}{d\theta}=\frac{F'(x)}{V''(x)} . \]
\[ \frac{dx_L}{d\lambda_S}=\frac{W'(x_L)}{V''(x_L)}>0,\qquad \frac{dx_L}{dK}=\frac{N'(x_L)}{V''(x_L)}<0,\qquad \frac{dy_{\text{post}}}{dK}=\frac{N'(y_{\text{post}})}{V''(y_{\text{post}})}>0 . \]
\[ \frac{d(y_{\text{post}}-x_L)}{dK}>0,\qquad N(x_L)=N(y_{\text{post}})+1 . \]
Control, $V''(x)=0$, $F'(x)=0$ $\Rightarrow$ $dx=0$ and $dx=1$ both solve the chain rule.
\end{mproof}

\begin{mproof}{M3}{Identification.det_formula, Identification.det_neg_of_opposite, Identification.det_neg, Identification.control_same_direction_can_fail}
\[ \det\begin{pmatrix}\frac{\rho_LW(y^*)}{\rho_L-\mu'}&\frac{\rho_LN(y^*)}{\rho_L-\mu'}\\[2pt] \frac{W'(x_L)}{V''(x_L)}&\frac{N'(x_L)}{V''(x_L)}\end{pmatrix}=\frac{\rho_L}{(\rho_L-\mu')V''(x_L)}\big(W(y^*)N'(x_L)-N(y^*)W'(x_L)\big) . \]
\[ a_{11},a_{12},a_{21}>0>a_{22}\ \Rightarrow\ a_{11}a_{22}<0<a_{12}a_{21}\ \Rightarrow\ a_{11}a_{22}-a_{12}a_{21}<0 . \]
Control, $a_{11}=a_{12}=a_{21}=a_{22}=1$ $\Rightarrow$ $\det=0$.
\end{mproof}

\begin{mproof}{M4}{SecondOrder.psi_identity, SecondOrder.psi_nonneg_of_qvi, SecondOrder.trigger_convex, SecondOrder.trigger_strict_iff, SecondOrder.control_trigger_flat_when_psi_zero}
\[ \tfrac12\sigma^2V''+\mu V'-\rho_LV-\Lambda=0,\quad V'=1+\kappa\ \Rightarrow\ \tfrac12\sigma^2V''(x_L^+)=\rho_LV+\Lambda-(1+\kappa)\mu=\Psi(x_L) . \]
\[ \mu(1+\kappa)-\rho_LV-\Lambda\le0\ \Rightarrow\ \Psi(x_L)\ge0\ \Rightarrow\ V''(x_L^+)\ge0,\qquad \sigma^2>0\ \Rightarrow\ \big(V''(x_L^+)>0\iff\Psi(x_L)>0\big) . \]
Control, $\sigma^2=1$, $V'=1$, $\kappa=0$, $V=\Lambda=\mu=0$, $\rho_L=1$ $\Rightarrow$ $\Psi=0$, $V''(x_L^+)=0$.
\end{mproof}

\begin{mproof}{M5}{SmoothFit.hasDeriv_Mv, SmoothFit.kink_is_convex, SmoothFit.no_convex_kink, SmoothFit.smooth_fit, SmoothFit.control_kink_needs_obstacle, Envelope.smf_permits}
\[ V\ge\mathcal MV,\ \ V(x_L)=\mathcal MV(x_L),\ \ (\mathcal MV)'=1+\kappa\ \Rightarrow\ V'(x_L^+)\ge1+\kappa . \]
\[ V'(x_L^+)>1+\kappa\ \Rightarrow\ \tfrac12\sigma^2q+(1+\kappa)\mu-\rho_LV-\Lambda\le0\ \ \forall q,\qquad q=\frac{2(\rho_LV+\Lambda-(1+\kappa)\mu+1)}{\sigma^2}\ \Rightarrow\ 1\le0 . \]
\[ \Rightarrow\ V'(x_L^+)=1+\kappa;\qquad q\le0,\ \ (1+\kappa)\mu\le\rho_LV+\Lambda\ \Rightarrow\ \tfrac12\sigma^2q+(1+\kappa)\mu-\rho_LV-\Lambda\le0 . \]
Control, $v(x)=x$, $\kappa=1$, $\mathcal Mv(x)=2x$ $\Rightarrow$ $v(0)=\mathcal Mv(0)$, $v'(0^+)=1<2$.
\end{mproof}

\begin{mproof}{M6}{ImpulseCount.partial_sums_bounded, ImpulseCount.count_summable, ImpulseCount.count_le, ImpulseCount.impulse_sum_summable, ImpulseCount.control_count_needs_fixed_cost}
\[ K\sum_{i<n}e_i\le B,\ K>0\ \Rightarrow\ \sum_{i<n}e_i\le\frac BK\ \ \forall n,\qquad e_i\ge0\ \Rightarrow\ N=\sum_ie_i\le\frac BK<\infty . \]
\[ |c_i|\le Ce_i\ \Rightarrow\ \sum_i|c_i|\le C\,N<\infty . \]
Control, $K=0$, $e_i=1$ $\Rightarrow$ $0\le0$ for every $n$ and $\sum_ie_i=\infty$.
\end{mproof}

\begin{mproof}{M7}{Targeting.target_free_of_start, Targeting.target_same_for_all_starts, Targeting.control_target_moves_with_convex_cost}
\[ V(z)-(1+\kappa)(z-x)-K=\big(V(z)-(1+\kappa)z\big)+\big((1+\kappa)x-K\big), \]
\[ \arg\max_z\big(V(z)-(1+\kappa)(z-x)-K\big)=\arg\max_z\big(V(z)-(1+\kappa)z\big)\quad\forall x . \]
Control, $V(z)=2z$, cost $(z-x)^2$ $\Rightarrow$ $\arg\max_z\big(2z-(z-x)^2\big)=x+1$.
\end{mproof}

\begin{mproof}{M8}{QVI_Part1.lambda_asymmetry_vanishes_at_one, QVI_Part1.neg_Lambda_concave, QVI_Part1.control_concavity_needs_lambda_ge_one}
\[ \Lambda(x)\big|_{\lambda_S=1}=(1-1)(R-x)^+=0 . \]
\[ -\Lambda(x)=(\lambda_S-1)\big(\min\{x,R\}-R\big),\qquad \lambda_S-1\ge0,\ \ x\mapsto\min\{x,R\}\ \text{concave} . \]
Control, $\lambda_S=0$, $R=0$ $\Rightarrow$ $-\Lambda(x)=\max\{-x,0\}$, $-\Lambda(0)=0<\tfrac12=\tfrac12(-\Lambda(-1))+\tfrac12(-\Lambda(1))$.
\end{mproof}

\begin{mproof}{M9}{CapGeometry.zmap_inverse, CapGeometry.zmap_of_inverse, Coordinate.zmap_strictMonoOn, Coordinate.zmap_at_requirement, Coordinate.zmap_tendsto_atBot, Coordinate.zmap_tendsto_atTop}
\[ z(x)=\log\frac{x-1}{q},\qquad 1+qe^{z(x)}=x\ \ (x>1),\qquad z(1+qe^{w})=w . \]
\[ 1<x<y\ \Rightarrow\ \frac{x-1}{q}<\frac{y-1}{q}\ \Rightarrow\ z(x)<z(y),\qquad z(1+q)=\log1=0 . \]
\[ \frac{x-1}{q}\to0^+\ (x\to1^+)\ \Rightarrow\ z\to-\infty,\qquad \frac{x-1}{q}\to\infty\ (x\to\infty)\ \Rightarrow\ z\to\infty . \]
\end{mproof}

\begin{mproof}{M10}{CapGeometry.hasDerivAt_zmap, CapGeometry.zvol2_eq, CapGeometry.xbar_sub_one, CapGeometry.one_lt_xbar, CapGeometry.solvency_binds, CapGeometry.liquidity_binds, CapGeometry.zvol2_solvency}
\[ \frac{x-1}{a_1}\le\frac{x-a_2}{a_3}\iff(a_3-a_1)\,x\le a_3-a_1a_2\iff x\le\bar x,\qquad \bar x-1=\frac{a_1(1-a_2)}{a_3-a_1}>0 . \]
\[ z'(x)=\frac{1}{x-1},\qquad \zeta(x)=\omega(x)^2\sigma^2-2\,\omega(x)\,c\,\sigma\sigma_L+\sigma_L^2 . \]
\[ 1<x\le\bar x,\ \pi(x)\,x=u(x)\ \Rightarrow\ \omega(x)=\frac{1}{a_1}\ \Rightarrow\ \zeta(x)=\nu_1^2. \]
\end{mproof}

\begin{mproof}{M11}{CapGeometry.deficit_limit, CapGeometry.surplus_limit, CapGeometry.variance_ratio_is_cap_geometry, CapGeometry.control_surplus_limit_needs_a1_lt_a3, CapGeometry.control_constancy_needs_binding}
Near $1^+$, $\bar x>1$ by M10, so $1<x<\bar x$ eventually and $\zeta(x)=\nu_1^2$ there.
For $x\ge\bar x$,
\[ \omega(x)=\frac{x-a_2}{a_3(x-1)}=\frac{1}{a_3}+\frac{1-a_2}{a_3}\,(x-1)^{-1}\to\frac{1}{a_3}, \]
\[ \zeta(x)=\omega(x)^2\sigma^2-2\,\omega(x)\,c\,\sigma\sigma_L+\sigma_L^2\to\nu_3^2 . \]
Controls, at $\sigma=1$, $\sigma_L=c=0$, $q=1$.
\[ a_1=\tfrac12,\ a_2=0,\ a_3=\tfrac14\ \Rightarrow\ u(x)=2(x-1),\ \zeta\equiv4\ne16=\nu_3^2. \]
\[ \pi\equiv0\ \Rightarrow\ \zeta(2)=0\ne4=\nu_1^2. \]
\end{mproof}

\begin{mproof}{M12}{Bounds.qv_le_max, Bounds.zvol2_eq_qv, Bounds.nu2_eq_qv, Bounds.qv_zero, Bounds.deficit_bound, Bounds.surplus_bound, Bounds.control_bound_needs_long_exposure, Bounds.control_bound_needs_compliance}
\[ \zeta(x)=Q(\omega),\qquad Q(\omega)=\omega^2\sigma^2-2\omega c\sigma\sigma_L+\sigma_L^2,\qquad \omega=\frac{\pi(x)\,x}{x-1}, \]
\[ Q(0)=\sigma_L^2,\qquad Q(1/a_i)=\nu_i^2 . \]
\[ \omega=tM,\ t\in[0,1]\ \Rightarrow\ Q(\omega)=(1-t)Q(0)+tQ(M)+\sigma^2M^2(t^2-t)\le\max\{Q(0),Q(M)\} . \]
\[ 0\le\pi,\ \ \pi x\le u(x)\le\frac{x-1}{a_1}\ \Rightarrow\ 0\le\omega\le\frac1{a_1}\ \Rightarrow\ \zeta(x)\le\max\{\sigma_L^2,\nu_1^2\} . \]
\[ \pi x\le\frac{x-a_2}{a_3}\ \Rightarrow\ \omega\le w(x)=\frac1{a_3}+\frac{1-a_2}{a_3(x-1)}\to\frac1{a_3},\qquad Q(w(x))\to\nu_3^2 . \]
Controls, $(\sigma,\sigma_L,c)=(1,1,1)$, $\omega=-1$ $\Rightarrow$ $Q=4>1=\max\{Q(0),Q(1)\}$; $(\sigma,\sigma_L,c)=(1,0,0)$, $\omega=2>1$ $\Rightarrow$ $Q=4>1$.
\end{mproof}

\section{Evidence and verification}
\label{app:evidence}

Every row of the model table is a theorem of Lean~4 with Mathlib. The Lean
column names the theorems, and the proof of each row in
Appendix~\ref{app:proofs} names the same ones. A theorem cited as
\path{Namespace.name} is declared in the file
\path{theory/lean/Namespace.lean}, a copy of the file of the same name under
\path{lean/mathlib/}. The development is built on the toolchain
\path{v4.32.0-rc1} against Mathlib commit \path{360da6f}. Every literature row
rests on a record under \path{lit/}, whose suite checks each quotation on its
page of the pinned source and builds and audits the record's Lean file. The
script \path{theory/verify.sh} runs the suites below on every build, and
\path{theory/VERIFICATION.md} records what they printed.

\begin{enumerate}
\item \emph{The Lean build and axiom audit.} Every file under
\path{theory/lean/} is a build root and builds. No source contains
\texttt{sorry}. Every declared theorem is audited with \texttt{\#print axioms},
and none depends on an axiom outside \texttt{propext},
\texttt{Classical.choice} and \texttt{Quot.sound}. Each file is identical to
its source in \path{lean/mathlib/}.
\item \emph{The skeleton checker.} Rules MS-1 to MS-9 of
\path{theory/checks/verify_manuscript.py}, each also run on a well-formed
sample it must accept and a faulty sample it must reject.
\item \emph{The literature records.} The suite of every record a literature
row cites.
\end{enumerate}

\paragraph{What Lean does not carry.} Buffer variance $\zeta$ is defined as
$z'(x)^2\sigma^2(x,\pi(x))$, and that this is the diffusion coefficient of
$z(X_t)$ is Itô's formula, which is not proved (A3). That the optimal policy
binds the cap near both ends of the state space is a hypothesis of M11, which
K2 and C2 state as their condition, supported only by the numerical solve of
the parent bundle.

M1 to M3 take as hypotheses, in raw form, the inputs of A6. That $V$ is the
value function is the verification theorem, not proved. Differentiability of
$V$ and of the thresholds in $\theta$ is the regularity lemma and the implicit
function theorem. The envelope identities $\partial_{\lambda_S}V=-W$ and
$\partial_KV=-N$ are Danskin's theorem. The interior equation enters
differentiated in $x$ and in $\theta$. The signs $W(y^*)>0$, $N(y^*)>0$,
$N'<0$ on the inaction region and $W'(x_L)>0$ come from hitting-time
arguments, and $W'(x_L)>0$ holds exactly when $\rho_LW(y_{\text{post}})<R-x_L$.
$V''(x_L)>0$ follows from M4 given $\Psi(x_L)>0$, the strict form of the
inequality the quasi-variational inequality gives on the intervention side.
$V''(y_{\text{post}})<0$ is a hypothesis, since its derivation in the parent
bundle fails (Appendix~\ref{app:refuted}). $\rho_L>\mu'(y^*)$ is a parameter
condition. K1 and C1 state these as their hypotheses.

M5 takes the viscosity supersolution property of $V$ at $x_L$ as a
hypothesis, in the form that a kink admits lower test functions of every
curvature. M6 takes the cost bound as a hypothesis, that each impulse costs at
least $K$ and the strategy's value is finite, with $e_i=e^{-\rho_L\tau_i}$ the
discounted impulse times and $c_i$ the impulse terms of the verification
argument. Interior regularity of $V$, $C^2$ on the inaction region and $C^3$
away from $R$ and $\bar x$, is assumed wherever a row differentiates $V$.

\paragraph{Illustrations.} No row rests on SymPy or on numerical sampling.

\section{Terminology}
\label{app:terms}

\footnotesize
\begin{longtable}{@{}P{4.0cm}P{2.5cm}P{9.5cm}@{}}
\caption{Terms, each bound to a symbol defined in Appendix~\ref{app:proofs}.}\label{tab:terms}\\
\toprule
Term & Symbol & Meaning \\
\midrule
\endfirsthead
\toprule
Term & Symbol & Meaning \\
\midrule
\endhead
Assets-to-liabilities ratio & $x$ & Total assets over total liabilities, the state of the model. \\
Required capital & $q$ & The capital regulation requires, as a ratio to liabilities. \\
Capital buffer & $z$ & The log of equity over required capital, $\log((x-1)/q)$ (M9). \\
Distress boundary & $x=1$ & The state at which equity is zero. \\
Regulatory cap & $u(x)$ & The largest risky exposure regulation allows, $\min\{(x-1)/a_1,(x-a_2)/a_3\}$. \\
Solvency branch & $(x-1)/a_1$ & The branch of the regulatory cap that binds for $x\le\bar x$ (M10). \\
Liquidity branch & $(x-a_2)/a_3$ & The branch of the regulatory cap that binds for $x\ge\bar x$ (M10). \\
Crossing point & $\bar x$ & The state at which the two branches are equal. \\
Cap geometry & $a_1,a_2,a_3$ & The parameters of the regulatory cap. \\
Asset volatility & $\sigma$ & The volatility of the risky asset. \\
Liability volatility & $\sigma_L$ & The volatility of the shock to liabilities. \\
Return-liability correlation & $c$ & The correlation between the risky return and the liability shock. \\
Asymmetry parameter & $\lambda_S$ & The weight on shortfall below the reference in the objective. \\
Shortfall penalty & $\Lambda$ & $(\lambda_S-1)(R-x)^+$, the running penalty below the reference $R$ (M8). \\
Buffer variance & $\zeta(x)$ & The squared diffusion coefficient of the capital buffer (A3). \\
Deficit limit & $\nu_1^2$ & The limit of buffer variance at the distress boundary (M11). \\
Surplus limit & $\nu_3^2$ & The limit of buffer variance deep in surplus (M11). \\
Fixed issuance cost & $K$ & The cost of an issue that does not scale with its size, per unit of liabilities (Appendix~\ref{app:scope}). \\
Proportional issuance cost & $\kappa$ & The cost of an issue per unit raised. \\
Dividend barrier & $y^*$ & The state at which dividends are paid. \\
Recapitalisation trigger & $x_L$ & The state at which capital is raised. \\
Injection target & $y_{\text{post}}$ & The state reached by an issue. \\
Expected discounted shortfall & $W$ & $\mathbb E_x\int_0^\infty e^{-\rho_Lt}(R-X_t)^+dt$ under the optimal policy, $-\partial_{\lambda_S}V$. \\
Discounted injection count & $N$ & $\mathbb E_x\sum_ie^{-\rho_L\tau_i}$ under the optimal policy, $-\partial_KV$. \\
Impulse operator & $\mathcal MV$ & The value of issuing at once to the best target, net of both issuance costs. \\
\bottomrule
\end{longtable}
\normalsize

\section{Primitives and scope}
\label{app:scope}

\paragraph{The balance sheet.} The bank holds a risky asset with volatility
$\sigma$ and a riskless asset, and owes liabilities with their own shock of
volatility $\sigma_L$ and correlation $c$ with the risky return (A1). The
riskless asset does not track the liabilities, unlike the reserve asset of
\citet{blackperold1992}, so the liability shock enters buffer variance
unhedged. Rows M9 to M12 depend on it.

\paragraph{The regulatory cap.} The capped exposure is
$u(x)=\min\{(x-1)/a_1,(x-a_2)/a_3\}$ with $0<a_1<a_3$ and $a_2<1$ (A2), the cap
of \citet{bayraktar2026}. Under these restrictions the solvency branch binds
near the distress boundary and the liquidity branch deep in surplus (M10).
With $a_1\ge a_3$ the surplus limit of M11 fails, by its control. With no
borrowing, $u(x)\le x$, the liquidity branch lies above the ceiling wherever
$a_2\le(1-a_3)x$ (\lean{BCVW.ceiling_tighter_than_lcr}, in the record of L2),
so the surplus limit of M11 is a statement about the leveraged case. Rows M10,
M11, M12, K2, C2 and C3 depend on it.

\paragraph{The fixed issuance cost.} $K$ is charged per unit of liabilities,
so the dollar cost of an issue is $K$ times liabilities (A5). This keeps the
reduction of the bank's problem to the single state $x$, which a fixed cost in
dollars would break (L4). M1 to M3, K1 and C1 depend on it.

\paragraph{The shortfall penalty.} The objective charges
$\Lambda(x)=(\lambda_S-1)(R-x)^+$ below the reference $R$ (A4). The model does
not say whose cost the penalty is. The reading adopted is that it stands for
the supervisor's restrictions below a capital threshold, which no row derives.
What $\lambda_S$ identifies in K1 is the weight of that penalty, under this
reading.

\paragraph{The policy.} That the optimal policy is a dividend barrier, a
recapitalisation trigger and an injection target, and that $V$ is the value
function, is the verification theorem of the parent bundle, taken as given
(A6). M1 to M5 depend on it.

\section{Refuted conjectures}
\label{app:refuted}

\paragraph{R1, strict concavity at the injection target.} The parent bundle
derives $V''(y_{\text{post}})<0$ from the sign change of
$\varphi'=V'-(1+\kappa)$ at $y_{\text{post}}$, positive below and negative above,
and from $\varphi\in C^3$, arguing that $\varphi''(y_{\text{post}})=0$ would give
a zero of even order. A zero of odd order can still have $\varphi''=0$. With
$\varphi(y)=-y^4/4$ and $y_{\text{post}}=0$, $\varphi'(y)=-y^3$ is positive for
$y<0$ and negative for $y>0$, while $\varphi''(0)=0$
(\lean{SecondOrder.refuted_target_strict}). The weak inequality survives and
strictness does not, so $V''(y_{\text{post}})<0$ enters M2 as a hypothesis
(Appendix~\ref{app:evidence}). C1 does not use it.

\section{Responses to the referee reports}

Not applicable. The paper has not been through review.

\section{Coverage of the referee comments}

Not applicable, for the same reason.

\end{document}
